\documentclass[a4paper,11pt]{article}

\usepackage[T1]{fontenc}
\usepackage[utf8]{inputenc}
\usepackage[margin=2.5cm]{geometry}
\usepackage{amsmath,amssymb}
\usepackage{xcolor}
\usepackage{tcolorbox}
\tcbuselibrary{skins,breakable}
\usepackage{enumitem}
\usepackage{booktabs}
\usepackage{array}
\usepackage{hyperref}
\usepackage{fancyhdr}
\usepackage{microtype}
\usepackage{lmodern}

% ---- Color palette ----
\definecolor{itbaBlue}{HTML}{3FA8F5}
\definecolor{itbaAmber}{HTML}{F2A32B}
\definecolor{itbaRed}{HTML}{F25C54}
\definecolor{itbaGreen}{HTML}{3DD68C}
\definecolor{darkBg}{HTML}{13192A}
\definecolor{lightGray}{HTML}{F5F6FA}
\definecolor{midGray}{HTML}{888888}

% ---- Hyperref setup ----
\hypersetup{
  colorlinks=true,
  linkcolor=itbaBlue,
  urlcolor=itbaBlue,
  pdftitle={Experimentos de Seguridad Criptografica},
  pdfauthor={ITBA -- Criptografia y Seguridad},
}

% ---- Header / Footer ----
\pagestyle{fancy}
\fancyhf{}
\renewcommand{\headrulewidth}{0.4pt}
\fancyhead[L]{\small\textcolor{midGray}{ITBA -- Criptografia y Seguridad}}
\fancyhead[R]{\small\textcolor{midGray}{Experimentos de Seguridad}}
\fancyfoot[C]{\small\textcolor{midGray}{\thepage}}

% ---- tcolorbox styles ----
% Main experiment box
\tcbset{
  experimentbox/.style={
    enhanced,
    breakable,
    colback=white,
    colframe=white,
    fonttitle=\bfseries\large,
    attach boxed title to top left={yshift=-3mm, xshift=5mm},
    boxed title style={
      colback=#1,
      colframe=#1,
      sharp corners,
      fontupper=\color{white}\bfseries\large,
    },
    top=8mm,
    left=4mm,
    right=4mm,
    bottom=4mm,
    sharp corners,
    drop shadow,
    borderline west={3pt}{0pt}{#1},
  }
}

% Memory hook box (amber left border)
\newtcolorbox{hookbox}[1][]{
  enhanced,
  breakable,
  colback=itbaAmber!8,
  colframe=itbaAmber,
  leftrule=4pt,
  rightrule=0pt,
  toprule=0pt,
  bottomrule=0pt,
  sharp corners,
  fonttitle=\bfseries\small\color{itbaAmber},
  title={Gancho Mnemonico},
  #1
}

% Attack box (red left border)
\newtcolorbox{attackbox}[1][]{
  enhanced,
  breakable,
  colback=itbaRed!7,
  colframe=itbaRed,
  leftrule=4pt,
  rightrule=0pt,
  toprule=0pt,
  bottomrule=0pt,
  sharp corners,
  fonttitle=\bfseries\small\color{itbaRed},
  title={Patron de Ataque},
  #1
}

% Key insight box (green left border)
\newtcolorbox{insightbox}[1][]{
  enhanced,
  breakable,
  colback=itbaGreen!8,
  colframe=itbaGreen,
  leftrule=4pt,
  rightrule=0pt,
  toprule=0pt,
  bottomrule=0pt,
  sharp corners,
  fonttitle=\bfseries\small\color{itbaGreen},
  title={Consecuencia Clave},
  #1
}

% Protocol step table styling
\newcolumntype{A}{>{\centering\arraybackslash\bfseries}p{3.5cm}}
\newcolumntype{B}{>{\centering\arraybackslash}p{1cm}}
\newcolumntype{C}{>{\arraybackslash}p{7.5cm}}

% ---- Document ----
\begin{document}

% ---- Title page ----
\begin{tcolorbox}[
  enhanced,
  colback=darkBg,
  colframe=darkBg,
  sharp corners,
  halign=center,
  valign=center,
  height=5cm,
]
  {\color{white}\Huge\bfseries Experimentos de Seguridad Criptografica}\\[8pt]
  {\color{itbaBlue}\large EAV \,$\cdot$\, CPA \,$\cdot$\, CCA \,$\cdot$\, Mac-Forge \,$\cdot$\, Sig-Forge}\\[12pt]
  {\color{midGray}\normalsize ITBA -- Criptografia y Seguridad \quad|\quad \today}
\end{tcolorbox}

\vspace{0.5cm}

% ---- TOC ----
\tableofcontents

\vspace{1cm}

% ================================================================
\section{Estructura General}
% ================================================================

Todos los experimentos siguen la misma estructura: el adversario $\mathcal{A}$ intenta distinguir cual de dos opciones eligio el sistema. Lo que cambia entre experimentos es \textbf{cuanto poder extra} tiene $\mathcal{A}$.

\begin{tcolorbox}[
  enhanced,
  colback=lightGray,
  colframe=itbaBlue,
  leftrule=3pt, rightrule=0pt, toprule=0pt, bottomrule=0pt,
  sharp corners,
]
\textbf{Jerarquia de seguridad:}
\[
  \underbrace{\text{EAV} \;\subseteq\; \text{Mult} \;\subseteq\; \text{CPA} \;\subseteq\; \text{CCA}}_{\text{confidencialidad}}
  \qquad\Big|\qquad
  \underbrace{\text{Mac-Forge} \;\longrightarrow\; \text{Sig-Forge}}_{\text{integridad}}
\]
\end{tcolorbox}

\vspace{0.3cm}
\noindent Cada flecha significa: ``si eres seguro contra $X$, eres seguro contra todo lo anterior''. La jerarquia va de \textit{menos poder del adversario} (EAV) a \textit{mas poder} (CCA).

% ================================================================
\section{PrivK\textsuperscript{eav} --- Eavesdropping}
% ================================================================

\begin{tcolorbox}[
  enhanced, breakable,
  colback=white, colframe=itbaBlue,
  borderline west={3pt}{0pt}{itbaBlue},
  sharp corners, drop shadow,
  top=4mm, left=4mm, right=4mm, bottom=4mm,
]
{\color{itbaBlue}\large\bfseries PrivK\textsuperscript{eav} --- Escucha Pasiva}\par\vspace{4pt}
\textbf{Pregunta:} ?`Puede $\mathcal{A}$, solo con \emph{ver} un criptograma, deducir cual de dos mensajes fue cifrado?
\end{tcolorbox}

\vspace{0.5cm}

\subsection{Protocolo}

\begin{center}
\begin{tabular}{ACB}
\toprule
\textbf{Adversario $\mathcal{A}$} & & \textbf{Challenger $\Pi$} \\
\midrule
Elige $m_0, m_1$ con $|m_0|=|m_1|$ & $\xrightarrow{m_0,\,m_1}$ & \\
 & & Genera $k \leftarrow \mathcal{K}$, $b \leftarrow \{0,1\}$ \\
 & $\xleftarrow{c}$ & Calcula $c = \mathrm{Enc}_k(m_b)$ \\
Emite $b' \in \{0,1\}$ & & \\
\bottomrule
\end{tabular}
\end{center}

\vspace{0.3cm}
$\mathcal{A}$ \textbf{gana} si $b' = b$. El esquema es EAV-seguro si:
\[
  \Pr[\text{PrivK}^{\mathrm{eav}}_{\mathcal{A},\Pi}(n) = 1] \leq \tfrac{1}{2} + \mathrm{negl}(n)
\]

\subsection{Descripcion informal}

$\mathcal{A}$ recibe \textbf{UN} criptograma y nada mas. El sistema es seguro si $\mathcal{A}$ no puede hacer mejor que tirar una moneda. Puede realizar cualquier computacion, ser tan inteligente como quiera --- simplemente no puede pedirle al sistema que cifre mensajes adicionales.

\begin{hookbox}
\textbf{``El soplon con la oreja en la puerta''} --- $\mathcal{A}$ intercepta una carta sellada que contiene uno de dos mensajes que el mismo eligio. Solo mirando el sobre, no puede saber cual mensaje esta adentro.
\end{hookbox}

% ================================================================
\section{PrivK\textsuperscript{CPA} --- Chosen Plaintext Attack}
% ================================================================

\begin{tcolorbox}[
  enhanced, breakable,
  colback=white, colframe=itbaGreen,
  borderline west={3pt}{0pt}{itbaGreen},
  sharp corners, drop shadow,
  top=4mm, left=4mm, right=4mm, bottom=4mm,
]
{\color{itbaGreen}\large\bfseries PrivK\textsuperscript{CPA} --- Ataque de Texto Plano Elegido}\par\vspace{4pt}
\textbf{Pregunta:} Si $\mathcal{A}$ puede cifrar cualquier mensaje que quiera, ?`puede adivinar cual de sus dos mensajes cifro el sistema?
\end{tcolorbox}

\vspace{0.5cm}

\subsection{Protocolo}

\begin{center}
\begin{tabular}{ACB}
\toprule
\textbf{Adversario $\mathcal{A}$} & & \textbf{Challenger $\Pi$} \\
\midrule
 & & Genera $k \leftarrow \mathcal{K}$ \\
Consultas al oraculo $\mathrm{Enc}_k(\cdot)$ (ilimitadas) & $\rightleftharpoons$ & \\
Elige $m_0, m_1$ con $|m_0|=|m_1|$ & $\xrightarrow{m_0,\,m_1}$ & \\
 & & Elige $b \leftarrow \{0,1\}$ \\
 & $\xleftarrow{c}$ & Calcula $c = \mathrm{Enc}_k(m_b)$ \\
Sigue usando el oraculo & $\rightleftharpoons$ & \\
Emite $b' \in \{0,1\}$ & & \\
\bottomrule
\end{tabular}
\end{center}

\vspace{0.3cm}
$\mathcal{A}$ \textbf{gana} si $b' = b$. El esquema es CPA-seguro si:
\[
  \Pr[\text{PrivK}^{\mathrm{cpa}}_{\mathcal{A},\Pi}(n) = 1] \leq \tfrac{1}{2} + \mathrm{negl}(n)
\]

\subsection{Consecuencia fundamental}

\begin{insightbox}
\textbf{CPA-Seguro $\Rightarrow$ el cifrado debe ser probabilistico.}\\[4pt]
Cualquier cifrado \emph{determinista} falla CPA automaticamente: $\mathcal{A}$ consulta $\mathrm{Enc}_k(m_0)$ y $\mathrm{Enc}_k(m_1)$, compara con $c$, y gana con probabilidad 1.
\end{insightbox}

\subsection{Ataque sobre esquema determinista}

\begin{attackbox}
\begin{enumerate}[leftmargin=*, label=\textbf{\arabic*.}]
  \item Consultar el oraculo: $c_0 = \mathrm{Enc}_k(m_0)$ y $c_1 = \mathrm{Enc}_k(m_1)$.
  \item Recibir desafio $c$.
  \item Si $c = c_0 \Rightarrow b'=0$;\quad si $c = c_1 \Rightarrow b'=1$.
  \item Como el cifrado es determinista, $\Pr[\text{ganar}] = 1$.
\end{enumerate}
\end{attackbox}

\begin{hookbox}
\textbf{``La fotocopiadora de mensajes secretos''} --- $\mathcal{A}$ tiene acceso ilimitado a la maquina de cifrado enemiga. El sistema debe resistir aun asi. La unica salida: producir un criptograma diferente cada vez (cifrado probabilistico).
\end{hookbox}

% ================================================================
\section{PrivK\textsuperscript{CCA} --- Chosen Ciphertext Attack}
% ================================================================

\begin{tcolorbox}[
  enhanced, breakable,
  colback=white, colframe=itbaRed,
  borderline west={3pt}{0pt}{itbaRed},
  sharp corners, drop shadow,
  top=4mm, left=4mm, right=4mm, bottom=4mm,
]
{\color{itbaRed}\large\bfseries PrivK\textsuperscript{CCA} --- Ataque de Texto Cifrado Elegido}\par\vspace{4pt}
\textbf{Pregunta:} Si $\mathcal{A}$ puede cifrar \emph{y} descifrar cualquier cosa (salvo el desafio exacto), ?`puede distinguir?
\end{tcolorbox}

\vspace{0.5cm}

\subsection{Protocolo}

\begin{center}
\begin{tabular}{ACB}
\toprule
\textbf{Adversario $\mathcal{A}$} & & \textbf{Challenger $\Pi$} \\
\midrule
 & & Genera $k \leftarrow \mathcal{K}$ \\
Oracles $\mathrm{Enc}_k(\cdot)$ y $\mathrm{Dec}_k(\cdot)$ & $\rightleftharpoons$ & \\
Elige $m_0, m_1$ & $\xrightarrow{m_0,\,m_1}$ & \\
 & & Elige $b \leftarrow \{0,1\}$ \\
 & $\xleftarrow{c}$ & Calcula $c = \mathrm{Enc}_k(m_b)$ \\
Sigue usando oracles (\emph{excepto} $\mathrm{Dec}_k(c)$) & $\rightleftharpoons$ & \\
Emite $b' \in \{0,1\}$ & & \\
\bottomrule
\end{tabular}
\end{center}

\vspace{0.3cm}
$\mathcal{A}$ \textbf{gana} si $b' = b$. \textbf{Restriccion}: $\mathcal{A}$ no puede consultar $\mathrm{Dec}_k(c)$ directamente.

\subsection{Resultado clave}

\begin{insightbox}
Ningun cifrado estandar por si solo es CCA-seguro. Se necesita \textbf{cifrado autenticado (cifrado + MAC)} para lograr CCA-seguridad. La \emph{maleabilidad} de los cifrados sin autenticacion permite ataques creativos.
\end{insightbox}

\subsection{Ataque sobre cifrado de flujo}

\begin{attackbox}
Sea $c = m_b \oplus k$ el desafio (cifrado de flujo):
\begin{enumerate}[leftmargin=*, label=\textbf{\arabic*.}]
  \item Recibir desafio $c = m_b \oplus k$.
  \item Construir $c' = c \oplus (m_0 \oplus m_1)$. Nota: $c' \neq c$, permitido.
  \item Consultar $\mathrm{Dec}_k(c') = m_b \oplus (m_0 \oplus m_1)$, que es $m_1$ si $b=0$ o $m_0$ si $b=1$.
  \item Comparar con $m_0, m_1$ para determinar $b$. Exito con $\Pr = 1$.
\end{enumerate}
\end{attackbox}

\begin{hookbox}
\textbf{``El agente doble con un limite''} --- $\mathcal{A}$ trabaja dentro de la organizacion enemiga, puede enviar y recibir mensajes cifrados. La unica regla: no puede pasar el mensaje interceptado directamente al descifrador. Debe ser creativo (modificarlo levemente). Sin MAC, esa creatividad es suficiente.
\end{hookbox}

% ================================================================
\section{Mac-Forge --- Infalsificabilidad de MAC}
% ================================================================

\begin{tcolorbox}[
  enhanced, breakable,
  colback=white, colframe=itbaAmber,
  borderline west={3pt}{0pt}{itbaAmber},
  sharp corners, drop shadow,
  top=4mm, left=4mm, right=4mm, bottom=4mm,
]
{\color{itbaAmber}\large\bfseries Mac-Forge --- Infalsificabilidad de MAC}\par\vspace{4pt}
\textbf{Pregunta:} Con acceso ilimitado al oraculo de etiquetado, ?`puede $\mathcal{A}$ producir un par (mensaje, etiqueta) valido para un mensaje que nunca consulto?
\end{tcolorbox}

\vspace{0.5cm}

\subsection{Protocolo}

\begin{center}
\begin{tabular}{ACB}
\toprule
\textbf{Adversario $\mathcal{A}$} & & \textbf{Challenger MAC} \\
\midrule
 & & Genera $k \leftarrow \mathcal{K}$ \\
Consultas al oraculo $\mathrm{Mac}_k(\cdot)$ & $\rightleftharpoons$ & \\
Sea $Q$ el conjunto de mensajes consultados & & \\
Emite $(m, t)$ con $m \notin Q$ & & \\
\bottomrule
\end{tabular}
\end{center}

\vspace{0.3cm}
$\mathcal{A}$ \textbf{gana} si $\mathrm{Vrfy}_k(m, t) = 1$ \textbf{y} $m \notin Q$.

\vspace{0.3cm}
\noindent\textbf{Diferencia clave con PrivK:} No hay bit $b$ que adivinar. El objetivo es \emph{falsificar}: producir una etiqueta valida para algo que nunca se consulto. La restriccion $m \notin Q$ evita el ataque trivial.

\subsection{Metodo general para romper un MAC inseguro}

\begin{enumerate}[leftmargin=*, label=\textbf{Paso \arabic*.}]
  \item \textbf{Encontrar debilidad estructural}: ?`Procesa el mensaje completo? ?`Es determinista? ?`Expone material de clave?
  \item \textbf{Consulta estrategica}: pedir etiqueta del mensaje ``especial'' (ceros, prefijo) para extraer informacion.
  \item \textbf{Construir falsificacion}: usar info extraida para construir $(m, t)$ con $m \notin Q$.
  \item \textbf{Verificar}: comprobar que $\mathrm{Vrfy}_k(m, t) = 1$.
\end{enumerate}

\subsection{Tres MACs rotos (patrones de examen)}

\begin{attackbox}[title={Patron de Ataque -- MAC 1: $\mathrm{Mac}_k(m) = G(k) \oplus m$}]
\begin{enumerate}[leftmargin=*, label=\textbf{\arabic*.}]
  \item Consultar $f(0\cdots0)$ para obtener $G(k) \oplus 0\cdots0 = G(k)$.
  \item Emitir $(1\cdots1,\; G(k) \oplus 1\cdots1)$.
  \item Verificacion: $G(k) \oplus (1\cdots1) = G(k) \oplus 1\cdots1$. \checkmark
\end{enumerate}
\end{attackbox}

\vspace{0.3cm}

\begin{attackbox}[title={Patron de Ataque -- MAC 2: $\mathrm{Mac}_k(m) = k \oplus \text{first}_\ell(m)$}]
\begin{enumerate}[leftmargin=*, label=\textbf{\arabic*.}]
  \item Consultar $f(0\cdots0)$ para obtener $k \oplus 0\cdots0 = k$.
  \item Ahora se puede etiquetar cualquier mensaje. Emitir $(m^*, k \oplus \text{first}_\ell(m^*))$ con $m^* \notin Q$.
  \item Alternativa: tomar cualquier $m$, agregar bits al final (no verificados por el MAC).
\end{enumerate}
\end{attackbox}

\vspace{0.3cm}

\begin{attackbox}[title={Patron de Ataque -- MAC 3: $\mathrm{Mac}_k(m) = \mathrm{Enc}_k(|m|)$}]
\begin{enumerate}[leftmargin=*, label=\textbf{\arabic*.}]
  \item El MAC solo verifica la \emph{longitud} del mensaje.
  \item Consultar $f(\texttt{"ab"})$ para obtener $t = \mathrm{Enc}_k(2)$.
  \item Emitir $(\texttt{"xy"}, t)$. Como $|\texttt{"xy"}| = |\texttt{"ab"}| = 2$, la verificacion pasa. \checkmark
\end{enumerate}
\end{attackbox}

\begin{hookbox}
\textbf{``El falsificador de documentos notariales''} --- $\mathcal{A}$ tiene un notario que firma cualquier cosa. Condicion de victoria: presentar un documento firmado que el notario nunca vio. Si el MAC tiene alguna regularidad (usa solo parte del mensaje, expone clave, es aditivo), $\mathcal{A}$ lo explota.
\end{hookbox}

% ================================================================
\section{Sig-Forge --- Infalsificabilidad de Firma Digital}
% ================================================================

\begin{tcolorbox}[
  enhanced, breakable,
  colback=white, colframe=darkBg,
  borderline west={3pt}{0pt}{darkBg},
  sharp corners, drop shadow,
  top=4mm, left=4mm, right=4mm, bottom=4mm,
]
{\color{darkBg}\large\bfseries Sig-Forge --- Infalsificabilidad de Firma Digital}\par\vspace{4pt}
\textbf{Pregunta:} Con $pk$ y acceso al oraculo de firma, ?`puede $\mathcal{A}$ generar una firma valida para un mensaje que nunca firmo?
\end{tcolorbox}

\vspace{0.5cm}

\subsection{Protocolo}

\begin{center}
\begin{tabular}{ACB}
\toprule
\textbf{Adversario $\mathcal{A}$} & & \textbf{Challenger Sig} \\
\midrule
 & & Genera $(sk, pk) \leftarrow \mathrm{Gen}$ \\
Recibe $pk$ & $\xleftarrow{pk}$ & \\
Consultas al oraculo $\mathrm{Sign}_{sk}(\cdot)$ & $\rightleftharpoons$ & \\
Sea $Q$ el conjunto de mensajes consultados & & \\
Emite $(m, \sigma)$ con $m \notin Q$ & & \\
\bottomrule
\end{tabular}
\end{center}

\vspace{0.3cm}
$\mathcal{A}$ \textbf{gana} si $\mathrm{Vrfy}_{pk}(m, \sigma) = 1$ \textbf{y} $m \notin Q$.

\subsection{Dos diferencias respecto a Mac-Forge}

\begin{enumerate}[leftmargin=*, label=\textbf{\arabic*.}]
  \item $\mathcal{A}$ tambien tiene $pk$ --- puede \emph{verificar sus propios intentos} sin consultar al oraculo.
  \item La verificacion usa clave publica --- cualquiera puede verificar, lo que otorga \textbf{no repudio}.
\end{enumerate}

\subsection{Ataques sobre RSA-Firma textbook}

\begin{attackbox}[title={Patron de Ataque -- RSA: Ataque sin mensajes (no-message attack)}]
\begin{enumerate}[leftmargin=*, label=\textbf{\arabic*.}]
  \item Elegir $\sigma$ al azar.
  \item Calcular $m = \sigma^e \bmod n$. Por definicion de RSA, $\sigma = m^d \bmod n$.
  \item El par $(m, \sigma)$ es una firma valida \emph{sin ninguna consulta al oraculo}.
\end{enumerate}
\end{attackbox}

\vspace{0.3cm}

\begin{attackbox}[title={Patron de Ataque -- RSA: Maleabilidad multiplicativa}]
Dado $s_1 = \mathrm{Sign}_{sk}(m_1) = m_1^d \bmod n$ y $s_2 = \mathrm{Sign}_{sk}(m_2) = m_2^d \bmod n$:
\begin{enumerate}[leftmargin=*, label=\textbf{\arabic*.}]
  \item Calcular $s = s_1 \cdot s_2 \bmod n = (m_1 \cdot m_2)^d \bmod n$.
  \item El par $(m_1 \cdot m_2 \bmod n,\; s)$ es una firma valida.
  \item Si $m_1 \cdot m_2 \bmod n \notin Q$, es una falsificacion exitosa.
\end{enumerate}
\textbf{Solucion:} Firmar $H(m)$ en lugar de $m$ (RSA con hash).
\end{attackbox}

\begin{hookbox}
\textbf{``Falsificar la firma de un funcionario''} --- $\mathcal{A}$ tiene $pk$ (puede verificar) y puede pedirle al funcionario que firme cualquier cosa. Victoria: presentar un documento firmado que el funcionario nunca vio. Ventaja sobre MAC: se pueden verificar propios intentos con $pk$ sin preguntar a nadie.
\end{hookbox}

% ================================================================
\section{Tabla Comparativa}
% ================================================================

\begin{center}
\renewcommand{\arraystretch}{1.5}
\begin{tabular}{>{\bfseries}p{2.8cm} p{2.5cm} p{2.5cm} p{2.5cm} p{2.5cm} p{2.5cm}}
\toprule
 & \textbf{EAV} & \textbf{CPA} & \textbf{CCA} & \textbf{Mac-Forge} & \textbf{Sig-Forge} \\
\midrule
Objetivo & Distinguir $m_0/m_1$ & Distinguir $m_0/m_1$ & Distinguir $m_0/m_1$ & Falsificar etiqueta & Falsificar firma \\
\addlinespace
Oracles disponibles & Ninguno & $\mathrm{Enc}(\cdot)$ & $\mathrm{Enc}(\cdot)$, $\mathrm{Dec}(\cdot)$ & $\mathrm{Mac}(\cdot)$ & $\mathrm{Sign}(\cdot)$ \\
\addlinespace
Restriccion & --- & --- & $\mathrm{Dec}(c) $ prohibido & $m \notin Q$ & $m \notin Q$ \\
\addlinespace
Clave publica & No & No & No & No & Si (pk) \\
\addlinespace
Necesita prob. & No (pero ayuda) & Si (obligatorio) & Si & No aplica & No aplica \\
\addlinespace
Construccion segura & OTP, PRG & PRF/PRP modo CTR & Cifrado autenticado & PRF con dominio extendido & RSA + Hash \\
\addlinespace
Propiedad garantizada & Confidencialidad basica & Confidencialidad CPA & Confidencialidad + integridad & Autenticidad & Autenticidad + no repudio \\
\bottomrule
\end{tabular}
\end{center}

\vspace{1cm}

\begin{tcolorbox}[
  enhanced,
  colback=darkBg,
  colframe=darkBg,
  sharp corners,
  halign=center,
]
{\color{itbaBlue}\bfseries Regla de oro:} {\color{white}Si necesitas confidencialidad \emph{y} autenticidad, usa \textbf{cifrado autenticado} (CCA-seguro). Nunca confundas confidencialidad con autenticidad: un mensaje cifrado sin MAC puede ser alterado sin saberlo.}
\end{tcolorbox}

\end{document}
